% Chapter 6: Panel DML and Rolling Window Methods
% ============================================================================
% This chapter extends DML to:
% - Panel data with fixed effects
% - Time-varying treatment effects via rolling windows
% ============================================================================

\chapter{Panel DML and Rolling Window Methods}
\label{ch:panel_dml}

\section{Introduction}\label{sec:ch06_introduction}

Chapter 5 introduced TemporalPLRDML for single time series with autocorrelation. Many economic applications involve \textbf{panel data}---repeated observations on multiple units over time. This chapter develops DML methods for panel data, addressing both fixed effects and time-varying treatment effects.

\keyresult{
\textbf{Chapter Objectives:}
\begin{itemize}
    \item Understand panel data structure and fixed effects
    \item Implement within-transformation for DML with fixed effects
    \item Derive cluster-robust standard errors for panel data
    \item Apply rolling window DML for time-varying treatment effects
    \item Handle structural breaks and regime changes
\end{itemize}
}

\subsection{The Panel Data Advantage}

Panel data combines cross-sectional and time series variation:
\begin{equation}
Y_{it} = \tau \cdot T_{it} + g(X_{it}) + \alpha_i + \gamma_t + \varepsilon_{it}
\end{equation}

where:
\begin{itemize}
    \item $i = 1, \ldots, N$ indexes individuals (e.g., states, firms, customers)
    \item $t = 1, \ldots, T$ indexes time periods
    \item $\alpha_i$ is an individual-specific \textbf{fixed effect}
    \item $\gamma_t$ is a time-specific \textbf{fixed effect}
\end{itemize}

\insight{
\textbf{Why Panel Data Strengthens Causal Inference}

Panel data allows us to:
\begin{enumerate}
    \item \textbf{Control for unobserved heterogeneity}: Fixed effects absorb time-invariant confounders
    \item \textbf{Exploit within-unit variation}: Compare each unit to itself over time
    \item \textbf{Increase precision}: More observations than pure cross-section or time series
    \item \textbf{Test parallel trends}: Verify identifying assumptions
\end{enumerate}
}

\subsection{Chapter Roadmap}

\begin{itemize}
    \item \textbf{Section 6.2}: Panel Data Fundamentals---fixed effects and within transformation
    \item \textbf{Section 6.3}: Panel DML---combining fixed effects with machine learning
    \item \textbf{Section 6.4}: Cluster-Robust Inference---accounting for within-unit correlation
    \item \textbf{Section 6.5}: Rolling Window DML---time-varying treatment effects
    \item \textbf{Section 6.6}: Applications
    \item \textbf{Section 6.7}: Exercises
\end{itemize}

% ============================================================================
\section{Panel Data Fundamentals}\label{sec:panel_fundamentals}
% ============================================================================

\subsection{Fixed Effects Model}

The fixed effects model assumes individual-specific intercepts:
\begin{equation}
Y_{it} = \tau \cdot T_{it} + X_{it}'\beta + \alpha_i + \varepsilon_{it}
\end{equation}

The key identification challenge: if $\alpha_i$ correlates with $T_{it}$ (e.g., more productive firms invest more), OLS on pooled data is biased.

\begin{definition}[Within Transformation]
\label{def:within_transform}
The \textbf{within transformation} demeans each variable by its individual-specific mean:
\begin{align}
\ddot{Y}_{it} &= Y_{it} - \bar{Y}_i \quad \text{where} \quad \bar{Y}_i = \frac{1}{T}\sum_{t=1}^{T} Y_{it} \\
\ddot{T}_{it} &= T_{it} - \bar{T}_i \\
\ddot{X}_{it} &= X_{it} - \bar{X}_i
\end{align}

This eliminates $\alpha_i$ since $\ddot{\alpha}_i = \alpha_i - \alpha_i = 0$.
\end{definition}

After within transformation:
\begin{equation}
\ddot{Y}_{it} = \tau \cdot \ddot{T}_{it} + \ddot{X}_{it}'\beta + \ddot{\varepsilon}_{it}
\end{equation}

OLS on the transformed data yields the \textbf{fixed effects estimator}.

\subsection{Two-Way Fixed Effects}

When both individual and time effects matter:
\begin{equation}
Y_{it} = \tau \cdot T_{it} + X_{it}'\beta + \alpha_i + \gamma_t + \varepsilon_{it}
\end{equation}

The two-way within transformation:
\begin{equation}
\ddot{Y}_{it} = Y_{it} - \bar{Y}_i - \bar{Y}_t + \bar{Y}
\end{equation}

where $\bar{Y}_t$ is the time-specific mean and $\bar{Y}$ is the grand mean.

\begin{remark}
Two-way fixed effects are particularly important for policy evaluation, where both cross-sectional heterogeneity (states differ) and common time shocks (recessions affect everyone) may confound treatment effects.
\end{remark}

% ============================================================================
\section{Panel DML}\label{sec:panel_dml}
% ============================================================================

\subsection{The Problem with Fixed Effects + ML}

Standard fixed effects require:
\begin{enumerate}
    \item Linear functional form: $g(X_{it}) = X_{it}'\beta$
    \item Known confounders: All relevant $X_{it}$ included
\end{enumerate}

DML relaxes the linearity assumption, but combining ML with fixed effects requires care.

\begin{algorithm}[H]
\caption{Panel DML with Fixed Effects}
\label{alg:panel_dml}
\begin{algorithmic}[1]
\REQUIRE Panel data $(Y_{it}, T_{it}, X_{it})_{i=1,t=1}^{N,T}$
\ENSURE Treatment effect $\hat{\tau}$ with cluster-robust SE

\STATE \textbf{Apply fixed effects transformation}:
\FOR{each variable $V \in \{Y, T, X\}$}
    \STATE Compute individual means: $\bar{V}_i = \frac{1}{T}\sum_t V_{it}$
    \STATE Demean: $\ddot{V}_{it} = V_{it} - \bar{V}_i$
\ENDFOR

\STATE \textbf{Run TemporalPLRDML on transformed data}:
\STATE $\hat{\tau} \gets \text{TemporalPLRDML}(\ddot{Y}, \ddot{T}, \ddot{X})$

\STATE \textbf{Compute cluster-robust standard errors}:
\FOR{each individual $i$}
    \STATE Compute cluster-level influence: $\Psi_i = \sum_t \psi_{it}$
\ENDFOR
\STATE Cluster variance: $\hat{V}_{\text{cluster}} = \frac{N}{N-1} \cdot \frac{1}{n^2}\sum_{i=1}^{N} \Psi_i^2$ \quad ($n$ = total retained observations)
\STATE Standard error: $\widehat{SE} = \sqrt{\hat{V}_{\text{cluster}}}$

\RETURN $\hat{\tau}$, $\widehat{SE}$
\end{algorithmic}
\end{algorithm}

\subsection{Implementation}

\begin{minted}[fontsize=\small]{python}
from dml_ts import PanelDML
import numpy as np

# Generate panel data
np.random.seed(42)
n_individuals = 100
n_periods = 20
n_total = n_individuals * n_periods

# Create IDs
individual_id = np.repeat(np.arange(n_individuals), n_periods)
time_id = np.tile(np.arange(n_periods), n_individuals)

# Individual fixed effects (unobserved)
alpha_i = np.random.randn(n_individuals)
alpha_expanded = alpha_i[individual_id]

# Covariates and treatment (correlated with fixed effects!)
X = np.random.randn(n_total, 3)
T = 0.5 * X[:, 0] + 0.3 * alpha_expanded + np.random.randn(n_total)

# Outcome with true effect tau = 2.0
true_tau = 2.0
Y = true_tau * T + X[:, 0]**2 + alpha_expanded + np.random.randn(n_total)

# Without fixed effects: BIASED
from dml_ts import TemporalPLRDML
naive_model = TemporalPLRDML(n_lags=0, model_y="ridge", model_t="ridge")
naive_result = naive_model.fit(Y, T, X)
print(f"Naive (no FE): theta = {naive_result.theta:.3f}")  # Biased!

# With fixed effects: UNBIASED
panel_model = PanelDML(
    fixed_effects="individual",
    cluster_se=True,
    model_y="ridge",
    model_t="ridge"
)
panel_result = panel_model.fit(Y, T, X, individual_id, time_id)
print(f"Panel DML:     theta = {panel_result.theta:.3f}")  # Close to 2.0
\end{minted}

\subsection{Monte Carlo Comparison}

\begin{table}[htbp]
\centering
\caption{Panel DML vs.\ Naive DML: Monte Carlo Results}
\label{tab:panel_monte_carlo}
\begin{tabular}{lccc}
\hline
\textbf{Method} & \textbf{Bias} & \textbf{Avg.\ SE} & \textbf{Coverage} \\
\hline
Naive DML (no FE) & $+0.48$ & 0.040 & 0\% \\
Panel DML (individual FE) & $+0.01$ & 0.046 & 90\% \\
Panel DML (two-way FE) & $+0.01$ & 0.046 & 90\% \\
\hline
\end{tabular}

\smallskip
\noindent\footnotesize Generated by
\texttt{scripts/mc\_panel\_dml\_table.py} (100 simulations, 40 units
$\times$ 15 periods, true $\tau = 2$, ridge nuisances, cluster-robust SEs).
The 90\% coverage at nominal 95\% reflects normal critical values with a
moderate cluster count (Remark above); $t_{N-1}$ critical values close most
of the gap.\normalsize
\end{table}

\insight{
\textbf{Key Finding}: Without fixed effects, DML has severe omitted variable bias ($+0.48$ on a true effect of $2.0$) and its confidence intervals never cover the truth. Panel DML with fixed effects removes the bias and restores near-nominal coverage.
}

% ============================================================================
\section{Cluster-Robust Inference}\label{sec:cluster_robust}
% ============================================================================

\subsection{Why Clustering Matters}

Within each individual, observations are typically correlated:
\[
\Cov(\varepsilon_{it}, \varepsilon_{is}) \neq 0 \quad \text{for } t \neq s
\]

This violates the i.i.d.\ assumption and inflates standard errors.

\begin{definition}[Cluster-Robust Variance]
\label{def:cluster_variance}
Since $\hat{\tau} - \tau_0 \approx \frac{1}{n}\sum_i \Psi_i$ with $n$ the
total number of retained observations, the cluster-robust (CR1) variance of
$\hat{\tau}$ with $N$ clusters is:
\begin{equation}
\hat{V}_{\text{CR}} = \frac{N}{N-1} \cdot \frac{1}{n^2} \sum_{i=1}^{N} \Psi_i^2
\end{equation}
where $\Psi_i = \sum_{t=1}^{T_i} \psi_{it}$ is the sum of influence scores
within cluster $i$. No centering term appears because
$\sum_i \Psi_i = \sum_{it} \psi_{it} = 0$ exactly, by the first-order
condition of the FWL estimator; $N$ counts only clusters that retain
observations after temporal-CV trimming.
\end{definition}

\begin{theorem}[Cluster-Robust Inference]
\label{thm:cluster_robust}
Under regularity conditions with $N \to \infty$ clusters:
\[
\frac{\hat{\tau} - \tau_0}{\widehat{SE}_{\text{CR}}} \convd N(0, 1)
\]
where $\widehat{SE}_{\text{CR}} = \sqrt{\hat{V}_{\text{CR}}}$ ($\hat{V}_{\text{CR}}$ already estimates the variance of $\hat{\tau}$; dividing by $N$ again is the sums-as-means scale error this companion once shipped).
\end{theorem}

\begin{remark}
The cluster-robust variance is valid whether or not within-cluster
correlation exists; with independent observations it converges to the
heteroskedasticity-only variance. Two practical caveats: with few clusters,
normal critical values are mildly anticonservative ($t_{N-1}$ critical values
are the standard refinement; Cameron \& Miller 2015), and clusters whose
observations are entirely dropped by temporal-CV trimming carry no
information — the companion excludes them from $N$ and warns.
\end{remark}

\subsection{Implementation Details}

\begin{minted}[fontsize=\small]{python}
def cluster_robust_se(influence_scores, cluster_ids):
    """
    Compute cluster-robust standard error.

    Parameters
    ----------
    influence_scores : array, shape (n,)
        Influence function values for each observation
    cluster_ids : array, shape (n,)
        Cluster identifiers

    Returns
    -------
    se : float
        Cluster-robust standard error
    """
    unique_clusters = np.unique(cluster_ids)
    N = len(unique_clusters)

    # Sum influence scores within each cluster
    cluster_sums = np.zeros(N)
    for i, cluster in enumerate(unique_clusters):
        mask = cluster_ids == cluster
        cluster_sums[i] = np.sum(influence_scores[mask])

    # CR1: V = N/(N-1) * sum(Psi_i^2) / n^2 — the sums are already
    # centered (sum of all influence scores is 0 by FWL construction).
    n = len(influence_scores)
    cluster_var = (N / (N - 1)) * np.sum(cluster_sums**2) / n**2

    return np.sqrt(cluster_var)
\end{minted}

% ============================================================================
\section{Rolling Window DML}\label{sec:ch06_rolling_window}
% ============================================================================

\subsection{Time-Varying Treatment Effects}

In many applications, treatment effects change over time:
\begin{itemize}
    \item Policy effects may fade or strengthen
    \item Market conditions evolve (e.g., competition intensity)
    \item Learning effects (customers adapt to prices)
\end{itemize}

Rolling window DML estimates local treatment effects:
\begin{equation}
\tau(t) = \text{ATE in window } [t - w/2, t + w/2]
\end{equation}

\begin{definition}[Rolling Window DML]
\label{def:rolling_window_dml}
For window size $w$ and step size $s$:
\begin{enumerate}
    \item For each center $t = w/2, w/2 + s, w/2 + 2s, \ldots$
    \item Select observations in $[t - w/2, t + w/2]$
    \item Estimate DML on window data
    \item Store $\hat{\tau}(t)$ and $\widehat{SE}(t)$
\end{enumerate}
\end{definition}

\subsection{Implementation}

\begin{minted}[fontsize=\small]{python}
from dml_ts import RollingWindowDML

# Fit rolling window DML
model = RollingWindowDML(
    window_size=100,   # 100 observations per window
    step_size=20,      # Move 20 observations between estimates
    model_y="ridge",
    model_t="ridge"
)

model.fit(Y, T, X, time_index=time)

# Get time series of local effects
time_centers, theta_series, se_series = model.get_effects()

# Plot time-varying effects
import matplotlib.pyplot as plt

plt.figure(figsize=(10, 5))
plt.plot(time_centers, theta_series, 'b-', label='Local ATE')
plt.fill_between(
    time_centers,
    theta_series - 1.96 * se_series,
    theta_series + 1.96 * se_series,
    alpha=0.3
)
plt.axhline(true_tau, color='r', linestyle='--', label='True ATE')
plt.xlabel('Time')
plt.ylabel('Treatment Effect')
plt.legend()
plt.title('Rolling Window DML: Time-Varying Treatment Effect')
plt.show()
\end{minted}

\subsection{Detecting Structural Breaks}

Rolling window DML naturally detects structural breaks---points where treatment effects change abruptly.

\begin{definition}[Structural Break Detection]
\label{def:structural_break}
A structural break at time $t^*$ is detected if:
\begin{equation}
\frac{|\hat{\tau}(t^* + s) - \hat{\tau}(t^* - s)|}{\sqrt{\widehat{SE}(t^*+s)^2 + \widehat{SE}(t^*-s)^2}} > z_{1-\alpha/2}
\end{equation}
where $z_{1-\alpha/2}$ is the critical value.
\end{definition}

\begin{remark}
For formal structural break testing, use the Chow test or Bai-Perron procedure. Rolling window DML provides visual diagnostics and approximate break detection.
\end{remark}

% ============================================================================
\section{Applications}\label{sec:applications}
% ============================================================================

\subsection{Application 1: Insurance Pricing}

\textbf{Setting}: $N = 50$ states, $T = 10$ years, estimating effect of competitor price changes on market share.

\textbf{Panel Structure}:
\begin{itemize}
    \item Individual FE: State-specific demand (regulation, demographics)
    \item Time FE: Macroeconomic conditions affecting all states
    \item Cluster SE: Observations within state are correlated
\end{itemize}

\textbf{Rolling Window}: Detect if price sensitivity changed after major regulatory reform.

\subsection{Application 2: Macroeconomic Policy}

\textbf{Setting}: Effect of monetary policy on output across $N = 30$ countries, $T = 40$ quarters.

\textbf{Panel Structure}:
\begin{itemize}
    \item Individual FE: Country-specific institutions
    \item Time FE: Global business cycle
    \item Two-way FE needed to avoid confounding
\end{itemize}

\textbf{Rolling Window}: Test if policy effectiveness changed after 2008 financial crisis.

% ============================================================================
\section{Recursive Dynamic G-Estimation}\label{sec:ch06_dynamic_g}
% ============================================================================

The estimators so far return a \emph{single} scalar effect. Many dynamic problems instead ask for a \emph{sequence} of period-specific effects: when a treatment is applied over $m$ periods and can alter both the outcome and the future state of the treated unit, what is the effect of the period-$\tau$ treatment? This is the dynamic treatment-effect setting of Lewis and Syrgkanis~\cite{lewisSyrgkanis2021}, and the companion class \texttt{DynamicGEstimationDML} implements their recursive g-estimation.

\keyresult{
\texttt{DynamicGEstimationDML} is the recursive dynamic g-estimator that the previous chapter's scalar \texttt{TemporalPLRDML} is \emph{not}: it recovers period-specific blips $\theta_1,\dots,\theta_m$ for a linear structural nested mean model, with valid inference. The constant (homogeneous) blip is implemented; heterogeneous blips $\theta_\tau(X)$ remain future work.
}

\subsection{Structural Nested Mean Models and the Blip}

Consider $n$ i.i.d.\ trajectories, each observed over $m$ periods, with treatment $T_\tau$, evolving controls (state) $X_\tau$, and a final outcome $Y$. Crucially, the state may depend on past treatment,
\[
  X_\tau = A X_{\tau-1} + b\,T_{\tau-1} + \varepsilon_\tau,
\]
so that past treatment confounds both future treatment and the outcome --- the \emph{time-varying confounding} that ordinary adjustment for $X$ mishandles (conditioning on a post-treatment state blocks part of the effect and opens collider paths). A \emph{linear structural nested mean model} writes the effect of the period-$\tau$ treatment on the final outcome as a constant \emph{blip} $\theta_\tau$:
\[
  Y = \sum_{\tau=1}^{m} \theta_\tau\, T_\tau + g(X) + U,
  \qquad \mathbb{E}[U \mid H_\tau, T_\tau] = \mathbb{E}[U \mid H_\tau],
\]
where $H_\tau = (X_{1:\tau}, T_{1:\tau-1})$ is the history available before $T_\tau$. Sequential conditional exchangeability given $H_\tau$ identifies $\theta_\tau$.

\subsection{Neyman-Orthogonal Recursive Peeling}

Lewis and Syrgkanis give a cross-fitted, Neyman-orthogonal g-estimator --- a dynamic generalization of Robinson's partially linear model --- that peels the blips off the outcome from the last period backward. Let $M_\tau$ denote the residual-maker that partials out $H_\tau$, with the conditional means $\mathbb{E}[\,\cdot \mid H_\tau]$ estimated by cross-fitting arbitrary machine-learning nuisances. For $\tau = m, m-1, \dots, 1$:
\begin{enumerate}
    \item \textbf{Residualize} on $H_\tau$: $\;\tilde{Y}^{(\tau)} = M_\tau\!\big(Y - \sum_{s>\tau} \theta_s T_s\big)$ and $\tilde{T}_\tau = M_\tau T_\tau$.
    \item \textbf{Solve} the orthogonal moment: $\;\theta_\tau = \langle \tilde{T}_\tau, \tilde{Y}^{(\tau)} \rangle / \langle \tilde{T}_\tau, \tilde{T}_\tau \rangle$.
    \item \textbf{Peel} period $\tau$ off the outcome and recurse to $\tau-1$.
\end{enumerate}
Inference uses the joint sandwich variance over the stacked moments. Because later-stage estimates enter earlier stages, Neyman orthogonality does \emph{not} hold \emph{across} the structural parameters: the Jacobian is triangular and the per-period variances are coupled. The implementation forms the full joint covariance rather than treating each period in isolation, with the sandwich ``meat'' clustered by unit (panel) or HAC/Newey--West (single series).

\subsection{Implementation}

\begin{minted}[fontsize=\small]{python}
from dml_ts import DynamicGEstimationDML
from dml_ts.validation import DynamicTreatmentDGP

# Known per-period blips with treatment-dependent state (time-varying confounding)
dgp = DynamicTreatmentDGP(
    n_periods=3, theta_t=[1.0, 0.5, 1.5], n_units=1500, p=3,
    state_feedback=True, treatment_state_coef=0.8, random_state=7,
)
data = dgp.generate()  # data.theta_t == [1.96, 2.90, 1.50] (total blips)

# Linear confounding here, so ridge nuisances are correctly specified
est = DynamicGEstimationDML(model_y="ridge", model_t="ridge", random_state=0)
result = est.fit(data.Y, data.T, data.X, groups=data.groups)
print(result.summary())
\end{minted}

The estimator recovers the true \emph{total} blips $(1.96, 2.90, 1.50)$ --- the direct coefficients plus the indirect path through the treatment-dependent state --- to within Monte Carlo error, whereas a naive pooled regression that ignores the sequential structure is badly biased. A reference path, \texttt{fit\_econml\_reference}, wraps EconML's \texttt{DynamicDML} (the reference Lewis--Syrgkanis implementation, optional via the \texttt{full} extra) and agrees with the native estimator to within sampling error --- a custom-versus-reference numerical cross-check.

\keyresult{
\textbf{Match the nuisance learner to the confounding.} Recursive peeling amplifies nuisance approximation error, so a flexible learner on a simple linear DGP can be \emph{worse} than a correctly specified linear one. Use \texttt{ridge} or linear nuisances for linear confounding, and \texttt{random\_forest} or \texttt{gradient\_boosting} only when the confounding is genuinely nonlinear.
}

A single long, stationary series is also supported (the \texttt{series} mode), where the blips become distributed-lag effects $Y_t = \sum_k \theta_{k+1}\, T_{t-k} + g(X_t) + U_t$ estimated with HAC inference.

% ============================================================================
\section{Exercises}\label{sec:ch06_exercises}
% ============================================================================

\begin{exercise}[Within Transformation]
\label{ex:within_transform}
Consider panel data with $N=3$ individuals and $T=2$ periods:

\begin{center}
\begin{tabular}{ccc|cc}
$i$ & $t$ & $Y_{it}$ & $T_{it}$ & $\alpha_i$ \\
\hline
1 & 1 & 5 & 1 & 2 \\
1 & 2 & 7 & 2 & 2 \\
2 & 1 & 3 & 1 & 1 \\
2 & 2 & 4 & 1 & 1 \\
3 & 1 & 8 & 2 & 3 \\
3 & 2 & 12 & 4 & 3 \\
\end{tabular}
\end{center}

\begin{enumerate}
    \item[(a)] Compute the within-transformed variables $\ddot{Y}_{it}$ and $\ddot{T}_{it}$
    \item[(b)] Estimate $\tau$ using OLS on the transformed data
    \item[(c)] What is the true $\tau$ in this DGP?
\end{enumerate}
\end{exercise}

\textbf{Solution to Exercise~\ref{ex:within_transform}:}

\begin{enumerate}
    \item[(a)] Individual means: $\bar{Y}_1 = 6$, $\bar{Y}_2 = 3.5$, $\bar{Y}_3 = 10$

    Within-transformed $Y$: $\ddot{Y} = (-1, 1, -0.5, 0.5, -2, 2)'$

    Similarly for $T$: $\bar{T}_1 = 1.5$, $\bar{T}_2 = 1$, $\bar{T}_3 = 3$

    $\ddot{T} = (-0.5, 0.5, 0, 0, -1, 1)'$

    \item[(b)] OLS on transformed: $\hat{\tau} = \frac{\sum \ddot{Y} \ddot{T}}{\sum \ddot{T}^2} = \frac{0.5 + 0.5 + 0 + 0 + 2 + 2}{0.25 + 0.25 + 0 + 0 + 1 + 1} = \frac{5}{2.5} = 2.0$

    \item[(c)] The true $\tau = 2.0$ (coefficient on $T$ in the DGP).
\end{enumerate}

\begin{exercise}[Cluster-Robust SE]
\label{ex:cluster_se}
You estimate $\hat{\tau} = 1.5$ from $n = 40$ observations in $N=4$
clusters, with cluster-level influence sums
\[
\Psi_1 = 0.2, \quad \Psi_2 = -0.1, \quad \Psi_3 = 0.3, \quad \Psi_4 = -0.4
\]
(they sum to zero, as FWL guarantees).

\begin{enumerate}
    \item[(a)] Compute the cluster variance $\hat{V}_{\text{CR}}$
    \item[(b)] Compute the cluster-robust standard error
    \item[(c)] Construct a 95\% confidence interval
\end{enumerate}
\end{exercise}

\textbf{Solution to Exercise~\ref{ex:cluster_se}:}

\begin{enumerate}
    \item[(a)] $\sum_i \Psi_i^2 = 0.04 + 0.01 + 0.09 + 0.16 = 0.30$

    $\hat{V}_{\text{CR}} = \frac{4}{3} \cdot \frac{0.30}{40^2} = \frac{0.40}{1600} = 0.00025$

    \item[(b)] $\widehat{SE} = \sqrt{0.00025} = 0.0158$

    \item[(c)] $1.5 \pm 1.96 \times 0.0158 = [1.469, 1.531]$
\end{enumerate}

% ============================================================================
\section{Summary}\label{sec:ch06_summary}
% ============================================================================

This chapter extended DML to panel data and time-varying effects:

\begin{enumerate}
    \item \textbf{Panel DML}: Within transformation eliminates fixed effects while preserving ML flexibility for nuisance functions

    \item \textbf{Cluster-Robust SE}: Accounting for within-cluster correlation prevents underestimated standard errors

    \item \textbf{Rolling Window DML}: Local estimation reveals how treatment effects change over time

    \item \textbf{Structural Breaks}: Rolling windows provide visual and approximate detection of regime changes
\end{enumerate}

\textbf{Key Takeaways:}
\begin{itemize}
    \item Always use fixed effects when unobserved heterogeneity correlates with treatment
    \item Cluster standard errors by the unit of treatment assignment
    \item Use rolling windows to test stability of treatment effects
    \item Two-way fixed effects handle both cross-sectional and time confounders
\end{itemize}

% ============================================================================
\section{Roadmap to Chapter 7}\label{sec:ch06_roadmap}
% ============================================================================

Chapter 7 integrates our methods with real macroeconomic data from FRED (Federal Reserve Economic Data). We develop:

\begin{enumerate}
    \item \textbf{FRED Data Loader}: Automated retrieval of macro indicators
    \item \textbf{Time Alignment}: Handling mixed-frequency data
    \item \textbf{Macro Control Variables}: GDP, CPI, interest rates as confounders
    \item \textbf{Application}: Policy effect estimation with real data
\end{enumerate}

This bridges the gap from synthetic validation to audited applications with real economic data.
